\scanpage{52}
(ii) Par une construction reposant \unsure{ess$^{\mathrm{t}}$.} sur les $\overline{\mathrm m}$ principes, mais sans doute un peu plus sophistiquée techniquement, (cf appendice), on trouve que si $f$ est lisse \unsure{qu. cpt} et qu. sép., alors
\[
\mathbb{R}f_*(\Omega^\bullet_{X/S})
\]
est muni d'une \uline{connexion absolue canonique} [au sens des catégories dérivées].
Cela ne signifie pas tt. à fait que les $\mathbb{R}^if_*(\Omega^\bullet_{X/S})$ sont munis individuellement de connexions absolues; c'est pourtant vrai si $S$ de car. 0, et $f$ propre et lisse (car alors la formation des $\mathbb{R}^i\unsure{f_*|f}(\Omega^\bullet_{X/S})$ eux-mêmes commute à l'extension de la base ....) --- ou encore, bien sûr, dans le cas lisse affine envisagé en (i)....
La vérification (pour $f$ propre, lisse) que cette connexion est compatible avec la construction transcendante, est en fait triviale.

\uline{Remarques}. On \unsure{pouvait|pourrait} espérer que ces connexions canoniques proviennent en fait de \uline{stratifications} canoniques, de façon que l'on aurait ainsi

\margin{51}
